\documentclass[11pt]{article}
\usepackage[utf8]{inputenc}\usepackage[T1]{fontenc}\usepackage[french]{babel}
\usepackage{amsmath,amssymb,booktabs}\usepackage[margin=2.2cm]{geometry}
\begin{document}
\section*{Élargissement thermique d'une raie de photoassociation}

\paragraph{Ce que fait la formule actuelle.}
\[
\Gamma_{\text{thermal}} = \frac{1}{\lambda}\sqrt{\frac{k_BT}{\mu}},\qquad \mu=\frac{m_{\text{at}}}{2}
\]
C'est la dispersion de la vitesse \emph{relative} des deux atomes, divisée par $\lambda$. Mais la vitesse relative ne donne pas d'effet Doppler :
le photon absorbé transfère son impulsion $\hbar k$ au \emph{centre de masse} de la paire, et le décalage Doppler vaut $\vec k\cdot\vec V$, où $\vec V$ est la vitesse du centre de masse.

\paragraph{Deux effets distincts pour une paire.}
\begin{enumerate}
\item \textbf{Doppler de la paire} (mouvement du centre de masse, masse totale $M=2m_{\text{at}}$) :
\[
\Gamma_{\text{D,paire}} = \frac{1}{\lambda}\sqrt{\frac{k_BT}{2m_{\text{at}}}}
= \frac{1}{\sqrt2}\,\Gamma_{\text{D,atome}},
\qquad \Gamma_{\text{D,atome}} = \frac{1}{\lambda}\sqrt{\frac{k_BT}{m_{\text{at}}}} .
\]
La paire, deux fois plus lourde, a un Doppler $\sqrt2$ fois plus \emph{petit} qu'un atome seul.
\item \textbf{Énergie de collision} (mouvement relatif, masse réduite $\mu$) : la résonance est déplacée de $E/h$ (terme $E/\hbar$ de l'équation de Zelevinsky), et $E$ suit la distribution thermique. La largeur associée est d'ordre
\[
\Gamma_E \sim \frac{k_BT}{h},
\]
avec un profil asymétrique (queue du côté de l'énergie de collision). C'est la contribution propre aux molécules.
\end{enumerate}

\paragraph{Valeurs numériques} ($^{88}$Sr, $\lambda=689$~nm, écarts-types pour les Doppler) :
\begin{center}
\begin{tabular}{lcccc}
\toprule
$T$ & $\Gamma_{\text{D,atome}}$ & $\Gamma_{\text{D,paire}}$ & formule actuelle ($\mu$) & $k_BT/h$ \\
\midrule
$1.4~\mu$K & $16.7$~kHz & $11.8$~kHz & $23.6$~kHz & $29$~kHz \\
$1.8~\mu$K & $18.9$~kHz & $13.4$~kHz & $26.8$~kHz & $37$~kHz \\
\bottomrule
\end{tabular}
\end{center}
La formule actuelle tombe par hasard entre les deux effets. Toutes ces largeurs restent petites devant les $120$--$340$~kHz mesurés : les conclusions ne changent pas.

\paragraph{Proposition de paragraphe pour la thèse.}
\begin{quote}\small
For molecular resonances, the thermal motion of the atoms broadens the line in two
ways. First, the Doppler shift is set by the motion of the centre of mass of the
colliding pair, which absorbs the photon recoil: since the pair has a mass
$2m_{\text{at}}$, its Doppler width $\frac{1}{\lambda}\sqrt{k_BT/2m_{\text{at}}}$ is $\sqrt2$
smaller than that of a single atom. Second, and specific to photoassociation, the
relative motion of the two atoms enters through the collision energy $E$, which shifts
the resonance by $E/h$: its thermal distribution broadens the line asymmetrically, by
an amount of order $k_BT/h$.
\end{quote}
\end{document}
